\chapter*{AGRADECIMIENTOS}
\phantomsection
\addcontentsline{toc}{chapter}{AGRADECIMIENTOS}
 
Al finalizar este trabajo, agradezco a quienes me acompañaron durante mi formación y contribuyeron, desde distintos espacios de mi vida, a que pudiera llegar hasta este momento.

A mi profesor guía, Patricio Galarce Acevedo, por su orientación, paciencia y disposición durante el desarrollo de este trabajo de titulación. Agradezco sus conocimientos, sus observaciones y el tiempo dedicado a acompañar este proceso, que me permitió aprender a abordar los desafíos con mayor rigor y perseverancia.

A mi familia, por el cariño y el apoyo que han sostenido mi formación. Gracias por confiar en mí, por escucharme y por acompañar tanto las alegrías como las dificultades de este recorrido. Este logro también guarda el valor de su presencia y de todos esos gestos cotidianos que, aunque no aparecen en estas páginas, hicieron una diferencia en mi camino.

Al Grupo Scout de Don Orione, por ser parte de mi formación humana y por las experiencias compartidas en torno al servicio, la fraternidad y el compromiso. A quienes me acompañaron en esa etapa, gracias por enseñarme el valor de colaborar, asumir responsabilidades y procurar que nuestras acciones aporten algo a los demás. Son aprendizajes que me acompañan más allá de cualquier espacio académico o profesional.

A mis compañeros del Club de Robótica, por compartir la curiosidad, el entusiasmo y las ganas de construir. Gracias por el intercambio de ideas, por aprender juntos y por hacer de la robótica un espacio de encuentro. En esa experiencia compartida encontré motivaciones que siguen presentes en mi manera de enfrentar nuevos proyectos.

A mis compañeros de la generación 2016 de Ingeniería Civil en Electrónica de la Universidad Tecnológica Metropolitana, por las jornadas de estudio, el intercambio de conocimientos y el compañerismo durante nuestra formación. Agradezco haber compartido con ustedes una etapa que dejó aprendizajes, recuerdos y vínculos que forman parte de mi historia.

A mis compañeros de PwC, por la generosidad al compartir sus conocimientos, por las conversaciones y por las oportunidades de aprender en el trabajo cotidiano. Gracias por el compañerismo y por contribuir a mi crecimiento profesional, ampliando mi forma de comprender los problemas y de buscar soluciones junto a otros.

A los docentes que fueron parte de mi formación y a quienes revisaron este trabajo, por sus enseñanzas, su tiempo y las observaciones que contribuyeron a mejorarlo.

Finalmente, agradezco a todas las personas que, con una palabra de ánimo, una idea o su compañía, estuvieron presentes durante este proceso. Al cerrar esta etapa, me llevo también la gratitud de haberla compartido con ustedes.